
Machine learning researchers frequently discover computational feasibility constraints only after full implementation, resulting in wasted resources on non-viable hypotheses. Existing approaches lack systematic early-stop protocols: ablation studies test hypothesis variations rather than viability, Big-O analysis omits constant factors, and expert intuition remains unvalidated. This work introduces Pilot-Driven Viability Gates, a framework treating viability assessment as incremental empirical validation across sample scales (10 samples, 100 samples, full dataset) with Bayesian posterior refinement. The central hypothesis is that computational overhead scales predictably from micro-pilot measurements, enabling accurate early prediction. Validation on a synthetic corpus of 32 machine learning hypotheses under a 10\% overhead threshold demonstrated 93.3\% accuracy (28 of 30 correct) at Gate 1 (10-sample micro-pilot) for identifying non-viable hypotheses, exceeding the 80\% target and significantly outperforming random baseline (binomial test p=4.$34\times 10^{-7})$. The framework exhibited perfect linear correlation (r=1.000, p<0.0001) between micro-pilot and full-scale overhead on synthetic data, with Bayesian updates reducing prediction error by 40.91\% (p=0.0003). These results establish proof-of-concept for the framework mechanics. External validity remains untested, as the synthetic corpus enforced perfect linearity by design. Real-world validation on published papers is necessary to determine whether the $r\geq 0.7$ correlation threshold holds under measurement noise and non-linear scaling effects.
